\documentclass[12pt]{article}

\usepackage[utf8]{inputenc}
\usepackage[T1]{fontenc}

\title{LA LINGÜÍSTICA DE LA BIOLOGÍA}
\author{E. CUELLAR, K. HERRERA, K. SANDOVAL, P. SANTIAGO}
\date{\today}

\begin{document}

\maketitle

\section{INTRODUCCIÓN}
Como muchos ya creen, el lenguaje no fue una invención biológica al azar, sino una adptación que surggió mediante la selección natural a lo largo de miles de años. Resulta necesario, por tanto, analizar el momento en que se reunieron las condiciones anatómicas, morfológicas y congnitivas de los \textbf{hominidos}. No debemos olvidar que no existió una única especie de homínidos, sino varias que convergieron evolutivamente y permitieron el surgimiento de la comunicación articulada.
El cambio más importante fue la adopción del bipedismo como forma de dezplazamiento. Caminar sobre dos extremidades no solo liberó nuestras manos para el empleamiento de herramientas, sino que también generó una reestructuracón de las vías respiratorias y del tracto vocal. El descenso de la laringe en la gargata creó una cámara de resonancia más amplia, lo que permitió al \textbf{\textit{Homo sapiens}} modular el flujo de aire y producir una enorme variedad de vocales y consonantes.
Sim embargo, el desarrollo de esta "cámara de resonancia" no fue el único factor evolutivo decisivo. Otros cambios resultaron igual de importantes, especialmente la encefalización progresiva, es decir, el desarrollo del cerebro y el espacio que este ocupa. En nuestra especie, este proceso no se limitó a un simple aumento del tamaño del craneo, sino que impulsó la especialización de áreas cerebrales altamente específicas. 

XX








... 


\section{BIBLIOGRAFÍA}
\begin{itemize}
  \item Hauser, M. D., Chomsky, N., \& Fitch, W. T. (2002). The faculty of language: What is it, who has it, and how did it evolve? \textit{Science}, \textit{298}(5598), 1569--1579. https://doi.org/10.1126/science.298.5598.1569
\end{itemize}

TBC (TO BE CONTINUED) 

\end{document}
